\chapter*{Mở đầu}
\addcontentsline{toc}{chapter}{Mở đầu}

Sự bùng nổ của thông tin số cùng nhu cầu cá nhân hóa ngày càng cao đã thúc đẩy sự chuyển dịch mạnh mẽ từ các hệ thống gợi ý tĩnh truyền thống sang các hệ thống gợi ý hội thoại (Conversational Recommender System - CRS). Bằng cách sử dụng ngôn ngữ tự nhiên làm phương thức tương tác chủ đạo, CRS cho phép hệ thống khai thác các sở thích phức tạp, linh hoạt và xử lý các nhu cầu mơ hồ của người dùng thông qua nhiều lượt đối thoại \cite{li2018towards, christakopoulou2016towardscrs}. Các bài toán lõi của CRS bao gồm việc nắm bắt chính xác ý định người dùng từ ngữ cảnh hội thoại, sinh ra các mục gợi ý phù hợp và tạo ra phản hồi ngôn ngữ tự nhiên để tương tác. Đây là một lĩnh vực thu hút sự quan tâm lớn từ cộng đồng nghiên cứu. Việc cá nhân hóa hội thoại đóng vai trò then chốt trong việc nâng cao trải nghiệm người dùng, giữ chân khách hàng và trực tiếp thúc đẩy giá trị thương mại cho các nền tảng dịch vụ.

Tuy nhiên, việc xây dựng một hệ thống CRS hoàn thiện đưa ra những thách thức đặc biệt. Sự xuất hiện của các Mô hình ngôn ngữ lớn (LLMs) đã cách mạng hóa năng lực xử lý ngôn ngữ, nhưng việc tích hợp LLM vào CRS hiện nay vẫn vấp phải rào cản cốt lõi: Tính tin cậy của thông tin (hiện tượng ảo giác) và giới hạn về năng lực suy luận đa bước. Khác với các hệ thống phân lớp hoặc dự đoán thông thường, hệ thống gợi ý hội thoại đòi hỏi hệ thống phải liên tục đối soát ý định của người dùng với cơ sở dữ liệu thực tế (như Đồ thị tri thức) để đảm bảo không gợi ý những mục không tồn tại.

Trong các nghiên cứu trước đây, nhiều kỹ thuật và kiến trúc khác nhau đã được đề xuất nhằm giải quyết bài toán này. Các hệ thống truyền thống tăng cường tri thức như KBRD \cite{chen-etal-2019-towards}, KGSF \cite{zhou2020improvingcrs} hay KERL \cite{Qiu_2025} đã tận dụng Đồ thị tri thức nhưng lại gặp hạn chế về khả năng hiểu ngôn ngữ linh hoạt. Các phương pháp dựa trên LLM theo hướng Zero-Shot (như LLM-ReDial \cite{liang-etal-2024-llm}) tỏ ra linh hoạt nhưng rất dễ sinh ra ảo giác do phụ thuộc hoàn toàn vào trí nhớ tham số. Các phương pháp tinh chỉnh như TallRec \cite{bao2023tallrec} hay ReFICR \cite{yang-2024-reficr} mang lại độ chính xác cao nhưng tốn kém tài nguyên và thiếu linh hoạt khi cập nhật dữ liệu mới. Gần đây, việc áp dụng kỹ thuật Truy xuất tăng cường sinh (Retrieval-Augmented Generation - RAG) trên Đồ thị tri thức đã cho thấy những kết quả đầy hứa hẹn. Điển hình, kiến trúc tiền đề GRACE (Graph-Reasoning Augmented Conversational Engine) đã giải quyết hiệu quả vấn đề ảo giác thông qua luồng truy xuất lai và sử dụng mô hình ngôn ngữ nhỏ để xếp hạng lại. Tuy nhiên, GRACE và các hệ thống RAG hiện tại vẫn hoạt động theo một "đường ống tĩnh". Khi đối mặt với các kịch bản hội thoại đa tầng, nơi các ràng buộc của người dùng mâu thuẫn nhau hoặc các truy vấn không có kết quả khớp tuyệt đối, kiến trúc tĩnh bộc lộ sự hạn chế: chúng thiếu đi cơ chế tự đánh giá, tự sửa lỗi và thỏa hiệp, dẫn đến việc trả về kết quả rỗng hoặc chất lượng thấp, làm gián đoạn mạch hội thoại.

Nhìn chung, các kiến trúc tĩnh hiện tại thiếu khả năng nhận thức tình huống và không thể phân rã, đối soát các yêu cầu phức tạp một cách độc lập. Để giải quyết những hạn chế đã nêu, khóa luận được thực hiện nhằm mục đích đề xuất hệ thống \textbf{ARGOS} (Adaptive Reasoning and Graph Orchestration System) chuyển đổi từ mô hình xử lý tuyến tính sang mô hình đa tác tử. Bằng cách phân chia các chức năng vào các tác tử chuyên biệt và thiết lập vòng lặp phản hồi, khóa luận hướng tới các đóng góp chính như sau:

(1) \textbf{Chuyển đổi kiến trúc từ tĩnh sang đa tác tử thích nghi:} Khóa luận đề xuất kiến trúc cộng tác giữa các tác tử: Orchestrator, Graph Agent, Critic Agent và Relaxation Agent. Khung làm việc này loại bỏ hoàn toàn điểm nghẽn của RAG dòng chảy tĩnh bằng cách đưa vào các vòng lặp phản hồi nội bộ. Hệ thống có khả năng tự động đánh giá lại chiến lược tìm kiếm, sử dụng cơ chế dự báo trọng số động để điều phối luồng truy xuất (Semantic, Content, Graph) dựa trên phân tích hình thái ngôn ngữ của truy vấn. Khóa luận xác thực tính hiệu quả của mô hình thông qua hai bộ dữ liệu hội thoại tiêu chuẩn là ReDial và INSPIRED.

(2) \textbf{Thiết lập cơ chế kiểm định nghiêm ngặt và triệt tiêu ảo giác:} Khóa luận trình bày một phương pháp đối soát chéo độc lập thông qua tác tử \textit{Critic Agent}. Bằng cách tách rời khâu trích xuất ràng buộc cứng và khâu kiểm định, mọi ứng viên được sinh ra đều phải trải qua bộ lọc của Critic Agent đối chiếu trực tiếp với Đồ thị tri thức Neo4j. Phương pháp này loại bỏ hoàn toàn các thực thể ảo giác do LLM tự sinh, đảm bảo độ chính xác và tính minh bạch tuyệt đối cho các gợi ý.

(3) \textbf{Khả năng tự sửa lỗi và thỏa hiệp thông minh:} Khi đối diện với các yêu cầu khắt khe dẫn đến tập kết quả rỗng, \textit{Relaxation Agent} được kích hoạt. Dựa trên thuật toán phân cấp ưu tiên, hệ thống có khả năng tự động phân tích điểm nghẽn, nới lỏng các ràng buộc ít quan trọng (ví dụ: năm phát hành) và giữ lại các giá trị cốt lõi (ví dụ: phong cách đạo diễn). Khả năng tự khởi động lại luồng truy vấn với ý định đã được tái cấu trúc giúp hệ thống khôi phục hội thoại một cách tự nhiên mà không cần đến sự can thiệp thủ công từ người dùng.

(4) \textbf{Tối ưu hóa độ trễ qua cơ chế Xếp hạng hai giai đoạn:} Nhằm đảm bảo khả năng đáp ứng thời gian thực, hệ thống đề xuất tách rời quá trình xếp hạng. Việc kết hợp mô hình Reranker thương mại nhẹ gọn cho vòng lọc đầu và giới hạn sự can thiệp của LLM lớn ở vòng suy luận cuối giúp giảm thiểu đáng kể chi phí tính toán và độ trễ, trong khi vẫn duy trì được hiệu suất gợi ý vượt trội. Thông qua cơ chế này, hệ thống thiết lập mốc hiệu năng SOTA mới so với các mô hình tinh chỉnh hàng đầu.

Phần còn lại của khóa luận được tổ chức như sau:

Chương 1 giới thiệu tổng quan, đặt vấn đề, tổng hợp các nghiên cứu liên quan, xác định khoảng trống nghiên cứu cùng mục tiêu và đóng góp của khóa luận. Chi tiết các kiến thức nền tảng về Hệ thống gợi ý hội thoại, Đồ thị tri thức, Mô hình ngôn ngữ lớn và kiến trúc Đa tác tử được trình bày tại Chương 2. Chương 3 đi sâu phân tích hệ thống tiền đề GRACE, chỉ ra những ưu điểm và định vị các điểm nghẽn của kiến trúc RAG tĩnh. Chi tiết các bước và cơ chế hoạt động của khung làm việc đa tác tử đề xuất ARGOS được mô tả kỹ lưỡng tại Chương 4. Chương 5 thể hiện kịch bản, thiết lập và các kết quả thực nghiệm, bao gồm phân tích chuyên sâu thông qua các chỉ số hiệu suất và so sánh với kết quả của các nghiên cứu hàng đầu hiện nay. Phần Kết luận được trình bày tại Chương 6, tóm lược các thành tựu chính đạt được và đề xuất các hướng phát triển trong tương lai.

Các nghiên cứu phân tích và cấu trúc mô hình tiền đề của khóa luận đã đóng góp vào nền tảng của bài báo khoa học\footnote{Thông tin bài báo: Phan Tiến Đạt, Trần Vỹ Anh, Hoàng Bảo Long, Trương Thị Minh Ngọc, Nguyễn Thị Hậu. "GRACE: A Knowledge Graph--Enhanced Conversational Recommendation System via Retrieval-Augmented Generation", Hội nghị Quốc tế về Công nghệ Thông tin và Truyền thông lần thứ 14 (SOICT 2025).} do người viết là tác giả chính \cite{phan2025grace}. Sự ghi nhận bước đầu về mặt học thuật này khẳng định tính đúng đắn của hướng nghiên cứu, đồng thời là cơ sở thực tiễn quan trọng để người viết tiếp tục phát triển và tối ưu hóa kiến trúc hệ thống đa tác tử trong phạm vi khóa luận này.




